% BEGIN EXACT CATALOG SOURCE: thm_10_10
\begin{thmbox}{10.10}
Let $(\Omega,\mathcal{F},P)$ be a probability space. For $n=1,2,3,\ldots$, let
\[
V_n(\omega)=(X_{n1}(\omega),X_{n2}(\omega),\ldots,X_{nd}(\omega))
\]
be a $d$-dimensional random vector, and let
\[
V(\omega)=(X_1(\omega),X_2(\omega),\ldots,X_d(\omega))
\]
be another $d$-dimensional random vector. Then:
\begin{enumerate}
\item $V_n\xrightarrow{\mathrm{a.s.}}V$ if and only if $X_{ni}\xrightarrow{\mathrm{a.s.}}X_i$ for all $i=1,\ldots,d$.
\item $V_n\xrightarrow{P}V$ if and only if $X_{ni}\xrightarrow{P}X_i$ for all $i=1,\ldots,d$.
\end{enumerate}
\end{thmbox}

\textit{Proof}
For part (a), suppose the component functions converge a.s. That is, for each $i=1,2,\ldots,d$, we can find an event $E_i$ with $P(E_i)=1$ such that
\[
X_{ni}(\omega)\xrightarrow{\mathrm{a.s.}}X_i(\omega)
\]
for $\omega\in E_i$. Then in $\bigcap_{i=1}^{d}E_i$, we have the convergence of random vector,
\[
\lim_{n\to\infty}V_n(\omega)=V(\omega)
\qquad\text{for }\omega\in\bigcap_{i=1}^{d}E_i.
\]
The proof is finished by noting that $P(\bigcap_{i=1}^{d}E_i)=1$.

Conversely, suppose $V_n\xrightarrow{\mathrm{a.s.}}V$. There is a set $E$ with probability $1$ such that $V_n(\omega)\to V(\omega)$ for $\omega\in E$. Each component of $V_n(\omega)$ converges to the corresponding component of $V$ for $\omega\in E$.

For part (b), assume that each component of $V_n$ converges in probability to the corresponding component of $V$. For any $\omega\in\Omega$,
\[
\lvert X_{ni}(\omega)-X_i(\omega)\rvert\leq \frac{\epsilon}{\sqrt{d}}
\text{ for all } i
\quad\Longrightarrow\quad
\lVert V_n(\omega)-V(\omega)\rVert\leq \epsilon.
\]
We thus have the following set inclusion:
\[
\{\omega:\lVert V_n-V\rVert>\epsilon\}
\subseteq
\bigcup_{i=1}^{d}
\left\{\omega:\lvert X_{ni}-X_i\rvert>\frac{\epsilon}{\sqrt{d}}\right\}.
\]
Taking probability of both sides, we get $P(\lVert V_n-V\rVert>\epsilon)\to 0$ as $n\to\infty$, which follows from the union bound (Exercise 2.4).

Conversely, assume that $V_n$ converges to $V$ in probability as a sequence of random vectors. Consider the $i$-th component of $V_n$. Since
\[
\lvert X_{ni}(\omega)-X_i(\omega)\rvert
\leq
\lVert V_n(\omega)-V(\omega)\rVert,
\]
we have
\[
\{\omega:\lvert X_{ni}(\omega)-X_i(\omega)\rvert>\epsilon\}
\subseteq
\{\omega:\lVert V_n(\omega)-V(\omega)\rVert>\epsilon\}.
\]
Hence, for any $\epsilon>0$,
\[
V_n\xrightarrow{P}V
\quad\Longrightarrow\quad
P(\lVert V_n(\omega)-V(\omega)\rVert>\epsilon)\to 0
\]
as $n\to\infty$, which implies
\[
P(\lvert X_{ni}(\omega)-X_i(\omega)\rvert>\epsilon)\to 0
\]
as $n\to\infty$. By the definition of convergence in probability, $X_{ni}$ converges to $X_i$ in probability for all $i$.
\hfill $\square$
% END EXACT CATALOG SOURCE: thm_10_10

% BEGIN EXACT CATALOG SOURCE: thm_10_11
\begin{thmbox}{10.11}
Let $(\Omega,\mathcal{F},P)$ be a probability space, let $V$ be a $d$-dimensional random vector, and let $(V_n)_{n=1}^{\infty}$ be a sequence of $d$-dimensional random vectors. Suppose $f:\mathbb{R}^d\to\mathbb{R}^m$ is a function that is continuous on a set $S_f\subseteq\mathbb{R}^d$ with $P(V\in S_f)=1$. Then:
\begin{enumerate}
\item $V_n\xrightarrow{\mathrm{a.s.}}V$ implies $f(V_n)\xrightarrow{\mathrm{a.s.}}f(V)$.
\item $V_n\xrightarrow{P}V$ implies $f(V_n)\xrightarrow{P}f(V)$.
\end{enumerate}
\end{thmbox}

\textit{Proof}
We let $f_k$ be the $k$-th component of the function $f$, for $k=1,2,\ldots,m$.

For part (a), by Theorem 10.10, it is sufficient to show that for each component function $f_k$ of $f$, we have
\[
f_k(V_n)\xrightarrow{\mathrm{a.s.}}f_k(V).
\]
From the hypotheses in the theorem, there exists a set $B$ with probability $1$ such that $V_n(\omega)\to V(\omega)$ for all $\omega\in B$. Let $A$ be the pre-image $V^{-1}(S_f)$. The set $A$ has probability $1$ by assumption. Hence, in $A\cap B$, we have
\[
\lim_{n\to\infty}V_n(\omega)=V(\omega),
\]
and
\[
\lim_{n\to\infty}f_k(V_n(\omega))=f_k(V(\omega)),
\]
by the continuity of $f_k$. Therefore $f_k(V_n)$ converges to $f_k(V)$ for all $\omega$ in $A\cap B$, which has probability $1$.

For part (b), by Theorem 10.10, it is sufficient to show that for each component function $f_k$ of $f$, we have $f_k(V_n)$ converging to $f_k(V)$ in probability. In the following we fix an index $k$.

Let $\delta$ and $\epsilon$ be positive real numbers. Consider the set
\[
A_{\delta,\epsilon}\coloneqq
\{v\in\mathbb{R}^d:\exists w\in\mathbb{R}^d
\text{ such that }\lVert w-v\rVert<\delta
\text{ and }
\lVert f_k(w)-f_k(v)\rVert>\epsilon\}.
\]
For each vector $v$ that is not in $A_{\delta,\epsilon}$, we cannot find any vector $w$ that satisfies $\lVert w-v\rVert<\delta$ and $\lVert f_k(w)-f_k(v)\rVert>\epsilon$ simultaneously. In set notation, we can express it as
\[
\{\omega:\lVert f_k(V_n(\omega))-f_k(V(\omega))\rVert>\epsilon\}
\subseteq
\{\omega:V(\omega)\in A_{\delta,\epsilon}\}
\cup
\{\omega:\lVert V_n(\omega)-V(\omega)\rVert\geq\delta\}.
\]
Take $\delta=1/j$ for positive integer $j$. By the continuity of $f$, and hence of the component function $f_k$, the set $A_{1/j,\epsilon}\cap S_f$ decreases to $\emptyset$ as $j$ tends to infinity. By the upper semi-continuous property of measure, we obtain
\[
P(\{\omega:V(\omega)\in A_{1/j,\epsilon}\})
=
P(\{\omega:V(\omega)\in A_{1/j,\epsilon}\cap S_f\})
\to 0
\]
as $j\to\infty$.

For an arbitrarily small $\epsilon_1>0$, we can find a sufficiently large $M$ such that
\[
P(\{\omega:V(\omega)\in A_{1/\ell,\epsilon}\})\leq \epsilon_1/2
\]
for all $\ell\geq M$.

Since it is assumed that $V_n\to V$ in probability, we can find a sufficiently large $N$ such that
\[
P(\{\lVert V_\ell(\omega)-V(\omega)\rVert\geq\delta\})\leq \epsilon_1/2
\]
for all $\ell\geq N$. This yields
\[
P(\{\lVert f_k(V_\ell(\omega))-f_k(V(\omega))\rVert>\epsilon\})
\leq \epsilon_1/2+\epsilon_1/2=\epsilon_1
\]
for all $\ell\geq \max(M,N)$. As $\epsilon_1$ is arbitrary, the probability
\[
P(\lVert f_k(V_\ell)-f_k(V)\rVert>\epsilon)
\]
converges to $0$ as $\ell\to\infty$.
\hfill $\square$
% END EXACT CATALOG SOURCE: thm_10_11
